\section{Вокруг языка C++}

\subsection{Какие стандарты языка существуют? Каков цикл выпуска новых стандартов?}
Язык C++ стандартизуется Международной организацией по стандартизации (ISO). Каждый стандарт фиксирует язык и стандартную библиотеку на некотором этапе развития.

\begin{definition}[Основные стандарты C++]
Исторически приняты следующие крупные ревизии языка:
\begin{itemize}
    \item \textbf{C++98} — первый международный стандарт;
    \item \textbf{C++03} — небольшое исправление и уточнение C++98;
    \item \textbf{C++11} — большой шаг с \textit{auto}, \textit{lambda}, \textit{move semantics}, \textit{range-based for};
    \item \textbf{C++14} — доработка C++11;
    \item \textbf{C++17} — развитие языка и библиотеки;
    \item \textbf{C++20} — крупное обновление: \textit{concepts}, \textit{modules}, \textit{ranges};
    \item \textbf{C++23} — очередная ревизия с улучшениями языка и библиотеки;
    \item \textbf{C++26} — следующий стандарт, который на момент обсуждения находится в стадии разработки как draft.
\end{itemize}
\end{definition}

\begin{remark}
В разговорной практике часто говорят не только о \textit{standard}, но и о \textit{draft} — рабочем проекте стандарта. Именно по draft обычно ориентируются implementaions компиляторов, когда реализуют новые возможности раньше окончательной публикации стандарта.
\end{remark}

После C++11 стандарты выпускаются примерно раз в три года: идея состоит в том, чтобы язык развивался регулярно, но без резких и слишком частых изменений.

\subsection{Что такое The Zero Overhead Principle? Привести пример появления этого принципа в коде.}
\begin{definition}[The Zero Overhead Principle]
\textit{Zero Overhead Principle} — принцип нулевых накладных расходов: если абстракция в C++ не используется, то она не должна ухудшать производительность по сравнению с низкоуровневым кодом; если используется, её цена должна быть минимальной и, по возможности, определяться на этапе компиляции.
\end{definition}

Идея языка состоит в том, чтобы давать высокоуровневые абстракции, но не заставлять платить за них лишнее время выполнения, память или дополнительные косвенные вызовы.

\begin{example}
Классический пример — шаблоны и `inline`-функции:

\begin{verbatim}
template <class T>
T max_value(const T& a, const T& b) {
    return (a < b) ? b : a;
}
\end{verbatim}

Здесь абстракция \textit{generic programming} устраняет необходимость писать отдельную функцию для каждого типа. После инстанцирования шаблона компилятор получает конкретный код для `int`, `double` и т.д., и при оптимизации он часто совпадает по стоимости с ручной реализацией.
\end{example}

Ещё один важный пример — RAII и умные указатели. Например, \texttt{std::unique\_ptr} даёт безопасное управление ресурсом, но в типичном случае не добавляет дорогой динамической диспетчеризации.

\subsection{Какие основные директивы языка CMake/файлов проекта? Для чего нужна директива \texttt{add\_executable}? Каковы её параметры?}
CMake — это не компилятор, а система генерации build system. В файле \texttt{CMakeLists.txt} описывают, что именно нужно собрать.

\begin{definition}[Типичные команды CMake]
На практике чаще всего встречаются:
\begin{itemize}
    \item \texttt{cmake\_min\_required(...)} — минимальная версия CMake;
    \item \texttt{project(...)} — имя проекта и языки;
    \item \texttt{add\_executable(...)} — создание исполняемого target;
    \item \texttt{add\_library(...)} — создание библиотеки;
    \item \texttt{target\_sources(...)} — привязка исходников к target;
    \item \texttt{target\_link\_libraries(...)} — линковка с библиотеками;
    \item \texttt{target\_include\_directories(...)} — пути к заголовкам;
    \item \texttt{set(...)} — установка переменных и свойств.
\end{itemize}
\end{definition}

\begin{definition}[\texttt{add\_executable}]
Команда \texttt{add\_executable} создаёт target типа \textit{executable}:
\[
\texttt{add\_executable(<name> [WIN32] [MACOSX\_BUNDLE] [EXCLUDE\_FROM\_ALL] <sources>...)}
\]
\end{definition}

Её смысл:
\begin{itemize}
    \item \texttt{<name>} — логическое имя цели сборки;
    \item \texttt{<sources>} — список исходных файлов;
    \item \texttt{WIN32} — пометка для Windows GUI-приложений;
    \item \texttt{MACOSX\_BUNDLE} — сборка в виде bundle на macOS;
    \item \texttt{EXCLUDE\_FROM\_ALL} — не включать цель в стандартную сборку по умолчанию.
\end{itemize}

\begin{example}
\begin{verbatim}
cmake_minimum_required(VERSION 3.20)
project(hello_cxx)

add_executable(hello main.cpp util.cpp)
\end{verbatim}

Здесь CMake создаёт исполняемый файл \texttt{hello} из двух исходников.
\end{example}

\subsection{Что такое интерфейс и что такое реализация? Что из них можно отнести к идеальному, а что — к материальному?}
\begin{definition}[Интерфейс]
\textit{Interface} — это то, \emph{что} объект, модуль или функция умеют делать и как с ними можно взаимодействовать. Интерфейс описывает контракт: сигнатуры функций, публичные типы, ограничения и обещания.
\end{definition}

\begin{definition}[Реализация]
\textit{Implementation} — это то, \emph{как} именно этот контракт выполняется внутри. Реализация содержит алгоритмы, данные, скрытые детали и технические решения.
\end{definition}

Интерфейс обычно относят к \textbf{идеальному}: он задаёт абстракцию, независимую от конкретного кода. Реализацию относят к \textbf{материальному}: это уже конкретный текст программы, который компилируется, линкуется и исполняется.

\begin{remark}
Хорошая архитектура старается минимизировать зависимость пользователя от реализации. Пользователь должен знать интерфейс, но не обязан знать внутреннее устройство.
\end{remark}

\subsection{Привести конкретные примеры разделения на интерфейс и реализацию в языке C++ (модули, типы и т.д.).}
Разделение на интерфейс и реализацию в C++ встречается на нескольких уровнях.

\begin{itemize}
    \item \textbf{Функции}: в \texttt{.h} объявление, в \texttt{.cpp} определение.
    \item \textbf{Классы}: в заголовке описывают публичные методы и поля, а в \texttt{.cpp} — их тела.
    \item \textbf{Абстрактные классы}: интерфейс задают pure virtual methods, а реализацию дают в наследниках.
    \item \textbf{Модули C++20}: \textit{module interface unit} экспортирует интерфейс, а \textit{module implementation unit} содержит скрытую реализацию.
    \item \textbf{PImpl idiom}: в заголовке виден только класс-обёртка, а настоящие поля спрятаны в реализации.
\end{itemize}

\begin{example}
\begin{verbatim}
// interface: shape.h
struct Shape {
    virtual double area() const = 0;
    virtual ~Shape() = default;
};

// implementation: circle.cpp
struct Circle : Shape {
    double r;
    double area() const override { return 3.14159 * r * r; }
};
\end{verbatim}

Здесь \texttt{Shape} — интерфейс, а \texttt{Circle} — конкретная реализация.
\end{example}

\subsection{Для чего используются заголовочные файлы в C++? Что в них обычно помещают? Что с ними обычно делают и что с ними обычно делать нельзя? Какие расширения встречаются у заголовочных файлов? Какая связь между заголовочными файлами и принципом интерфейса-реализации?}
\begin{definition}[Заголовочный файл]
\textit{Header file} — файл, который обычно подключают через \texttt{\#include} и который предназначен для распространения интерфейса между несколькими \textit{translation units}.
\end{definition}

Обычно в заголовках размещают:
\begin{itemize}
    \item объявления функций;
    \item определения классов и структур;
    \item объявления типов, alias-ов, enum-ов;
    \item \texttt{constexpr}, \texttt{inline}, шаблоны;
    \item константы и объявление внешних объектов, если они действительно должны быть общими.
\end{itemize}

Заголовок обычно:
\begin{itemize}
    \item \textbf{подключают} в \texttt{.cpp}-файлы;
    \item \textbf{не компилируют} как отдельную единицу;
    \item \textbf{защищают} include guard-ами или \texttt{\#pragma once}.
\end{itemize}

Что обычно \textbf{нельзя} делать в заголовке без особой необходимости:
\begin{itemize}
    \item размещать обычные определения функций и переменных, если это приведёт к множественным определениям;
    \item включать тяжёлые зависимости без нужды;
    \item создавать циклические зависимости между заголовками;
    \item путать интерфейс с внутренними деталями реализации.
\end{itemize}

Типичные расширения: \texttt{.h}, \texttt{.hpp}, \texttt{.hh}, \texttt{.hxx}, иногда \texttt{.inl} для inline-частей.

\begin{remark}
Связь с принципом interface/implementation здесь прямая: заголовок — это, как правило, место для интерфейса, а \texttt{.cpp}-файл — место для реализации.
\end{remark}

\subsection{Для чего используются модули трансляции в C++? Что в них обычно помещают? Что с ними обычно делают и что с ними обычно делать нельзя? Какие расширения встречаются у модулей трансляции? Какая связь между модулями трансляции и принципом интерфейса-реализации?}
\begin{definition}[Модульная трансляционная единица]
\textit{Module translation unit} — это translation unit, в котором есть \texttt{module declaration}. Такая единица может быть \textit{module interface unit}, \textit{module implementation unit} или \textit{partition}.
\end{definition}

Модули в C++20+ введены для более строгого и явного разделения интерфейса и реализации, чем в системе заголовков. Они уменьшают зависимость от текстового включения и помогают ускорять сборку.

Обычно в \textbf{module interface unit} помещают:
\begin{itemize}
    \item \texttt{export module ...;}
    \item экспортируемые объявления классов, функций, шаблонов;
    \item публичную часть API.
\end{itemize}

Обычно в \textbf{module implementation unit} помещают:
\begin{itemize}
    \item неэкспортируемые определения;
    \item скрытые helper-функции;
    \item детали реализации, не нужные пользователю модуля.
\end{itemize}

Что с модулями обычно \textbf{делают}:
\begin{itemize}
    \item импортируют через \texttt{import};
    \item экспортируют только нужную часть API;
    \item разделяют публичное и скрытое содержимое более строго, чем в headers.
\end{itemize}

Что с ними обычно \textbf{не делают}:
\begin{itemize}
    \item не используют как обычный заголовок с бесконтрольным текстовым \texttt{\#include};
    \item не дублируют экспорт без необходимости;
    \item не смешивают интерфейс и реализацию хаотично.
\end{itemize}

Типичные расширения: \texttt{.ixx}, \texttt{.cppm}, \texttt{.mpp}, иногда обычный \texttt{.cpp}.

\begin{remark}
Модули лучше выражают принцип interface/implementation: интерфейс задаётся через \texttt{export}, а реализация скрывается внутри модуля. В этом смысле модуль — более современный и формальный механизм, чем заголовочный файл.
\end{remark}
